% Un G-Counter sur trois répliques : chaque réplique incrémente sa propre case,
% la fusion prend le maximum case par case ; un message dupliqué ne change rien.
\documentclass[border=6pt]{standalone}
\usepackage{sd}
\begin{document}
\begin{tikzpicture}[x=1mm, y=1mm, note/.append style={font=\sffamily\small},
    vec/.style={draw=#1, fill=#1Bg, line width=0.7pt, rounded corners=2pt, inner sep=2pt,
      font=\ttfamily\small},
    lab/.style={font=\sffamily\small, text=#1, align=center},
    lab/.default=Dink]
  \foreach \r/\y/\c in {R1/0/Dblue, R2/-22/Dteal, R3/-44/Dpurple} {
    \node[lab=\c, font=\sffamily\small\bfseries, anchor=east] at (-4,\y) {\r};
    \draw[\c, line width=0.8pt, dash pattern=on 2pt off 2pt] (0,\y) -- (170,\y);
  }
  % incréments locaux
  \node[vec=Dblue] (a1) at (14,0) {[3, 0, 0]};
  \node[lab=Dslate, above=1pt] at (a1.north) {+3 sur R1};
  \node[vec=Dteal] (b1) at (14,-22) {[0, 5, 0]};
  \node[lab=Dslate, above=1pt] at (b1.north) {+5 sur R2};
  \node[vec=Dpurple] (c1) at (14,-44) {[0, 0, 2]};
  \node[lab=Dslate, below=1pt] at (c1.south) {+2 sur R3};
  % R1 -> R2
  \node[vec=Dteal] (b2) at (62,-22) {[3, 5, 0]};
  \draw[arr=Dblue] (a1.east) to[out=0, in=150] node[lab=Dblue, above, sloped, pos=0.5] {état de R1} (b2.north west);
  % R2 -> R3, puis doublon
  \node[vec=Dpurple] (c2) at (108,-44) {[3, 5, 2]};
  \draw[arr=Dteal] (b2.south east) to[out=-30, in=170] node[lab=Dteal, below, sloped, pos=0.5] {état de R2} (c2.west);
  \node[vec=Dpurple] (c3) at (150,-44) {[3, 5, 2]};
  \draw[arr=Dteal, dash pattern=on 3pt off 2pt] (b2.east) to[out=-10, in=150] node[lab=Dteal, above, sloped, pos=0.82] {en double} (c3.north west);
  % R3 -> R1
  \node[vec=Dblue] (a2) at (150,0) {[3, 5, 2]};
  \draw[arr=Dpurple] (c2.north) to[out=80, in=200] node[lab=Dpurple, right=3pt, pos=0.6] {état de R3} (a2.west);
  \node[lab=Dslate, anchor=north west, align=left] at (0,-54) {fusion : maximum case par case, donc commutative, associative et idempotente ;
    valeur = somme des cases = 10 partout ;\\le doublon ne change rien, alors qu’une fusion par addition compterait 5 en trop};
\end{tikzpicture}
\end{document}
